% year: תשע"ח
% type: מבחן
% semester: ב
% moed: א
% number: 2א
% points: 15
% categories: func-analysis, ivt
% source: exams/מבחנים/תשעח/מועד א תשעח סמסטר ב.pdf

\begin{question}
לאילו ערכים של פרמטר $a$ למשוואה $x^4-a\cdot x^3+27=0$ אין שורשים ממשיים?
\end{question}

\begin{hint}
$x=0$ אינו שורש. בודדו את הפרמטר: $a=\frac{x^4+27}{x^3}=x+\frac{27}{x^3}$, ומצאו את טווח הערכים של הפונקציה $g(x)=x+\frac{27}{x^3}$.
\end{hint}

\begin{solution}
$x=0$ אינו שורש (מתקבל $27\neq0$). עבור $x\neq0$:
\[
x^4-ax^3+27=0\iff a=\frac{x^4+27}{x^3}=x+\frac{27}{x^3}=:g(x).
\]
לכן למשוואה יש שורש ממשי אם ורק אם $a$ שייך לתמונה של $g$ על $\R\setminus\{0\}$. נמצא את התמונה.

\textbf{עבור $x>0$:} $g'(x)=1-\frac{81}{x^4}$, שמתאפסת ב-$x=3$; $g'<0$ ב-$(0,3)$ ו-$g'>0$ ב-$(3,\infty)$. לכן $x=3$ מינימום מוחלט על $(0,\infty)$, עם $g(3)=3+1=4$. בנוסף $g(x)\to+\infty$ כאשר $x\to0^+$ וכאשר $x\to+\infty$. מכיוון ש-$g$ רציפה, לפי משפט ערך הביניים התמונה של $(0,\infty)$ היא $[4,\infty)$.

\textbf{עבור $x<0$:} $g$ פונקציה אי-זוגית ($g(-x)=-g(x)$), ולכן התמונה של $(-\infty,0)$ היא $(-\infty,-4]$.

לכן התמונה של $g$ היא $(-\infty,-4]\cup[4,\infty)$, ולמשוואה יש שורש ממשי אם ורק אם $|a|\ge4$. (בדיקה: עבור $a=4$, $x=3$: $81-108+27=0$.)

\textbf{תשובה:} אין שורשים ממשיים אם ורק אם $-4<a<4$.
\end{solution}
